\documentclass[11pt,a4paper]{article}
\usepackage[utf8]{inputenc}
\usepackage[T1]{fontenc}
\usepackage{lmodern}
\usepackage{amsmath,amssymb,amsfonts,amsthm}
\usepackage{geometry}
\geometry{margin=1in}
\usepackage{xcolor}
\usepackage{hyperref}
\usepackage{cleveref}
\usepackage{booktabs}
\usepackage{array}
\usepackage{cite}
\usepackage{microtype}
\usepackage{tcolorbox}

\hypersetup{
    colorlinks=true,
    linkcolor=blue,
    citecolor=red,
    urlcolor=cyan,
    pdftitle={Pre-Registered Confrontation of a K3 x T2 Dual-Scale Hypothesis with Data: Four Observables, Four Negative Verdicts, and a Common Cause},
    pdfauthor={Xavier Callens and the SocrateAI Scientific Agora Collaboration}
}

\newtheorem{theorem}{Theorem}[section]
\newtheorem{proposition}[theorem]{Proposition}
\newtheorem{definition}[theorem]{Definition}
\theoremstyle{remark}
\newtheorem{remark}[theorem]{Remark}

% Epistemic tier badges, following papers 7--10.
\newcommand{\tierA}{\colorbox{green!15}{\textsf{\scriptsize Tier A --- Lean kernel}}}
\newcommand{\tierL}{\colorbox{blue!10}{\textsf{\scriptsize Tier L --- literature}}}
\newcommand{\tierC}{\colorbox{orange!15}{\textsf{\scriptsize Tier C --- conjecture}}}
\newcommand{\lean}[1]{\texttt{#1}}

\newtcolorbox{headlinebox}{%
  colback=red!4, colframe=red!45!black, boxrule=0.5pt, arc=1pt,
  left=6pt, right=6pt, top=4pt, bottom=4pt,
  fonttitle=\bfseries\small, title={Headline result}
}
\newtcolorbox{protocolbox}{%
  colback=gray!5, colframe=gray!45, boxrule=0.4pt, arc=1pt,
  left=6pt, right=6pt, top=3pt, bottom=3pt,
  fonttitle=\bfseries\small, title={Pre-registration protocol used in this paper}
}
\newtcolorbox{scopebox}{%
  colback=gray!5, colframe=gray!45, boxrule=0.4pt, arc=1pt,
  left=6pt, right=6pt, top=3pt, bottom=3pt,
  fonttitle=\bfseries\small, title={What is not proved}
}

\title{\Huge \textbf{Pre-Registered Confrontation of a}\\[0.3em]
\LARGE \textbf{$K3\times T^2$ Dual-Scale Hypothesis with Data:}\\[0.2em]
Four Observables, Four Negative Verdicts, and a Common Cause}

\author{\textbf{Xavier Callens}$^{1}$ \and \textbf{SocrateAI Scientific Agora Collaboration}$^{1}$ \\[0.5em]
$^{1}$\textit{SocrateAI Research Laboratory for Mathematical Physics \& Formal Verification}}

\date{September 2026}

\makeatletter\let\LMoldunderscore\_\renewcommand{\_}{\LMoldunderscore\allowbreak}\makeatother
\setlength{\emergencystretch}{3em}

\begin{document}

\maketitle

\begin{abstract}
A research programme built around a ``micro/macro dual-scale'' reading of string theory on
$K3\times T^2$ identifies the string scale with the geometric mean of the Planck length and the Hubble
radius, $\alpha'=\ell_P\,c/H_0$, and hence a self-dual length $s=\sqrt{\ell_P\,c/H_0}\approx47\,\mu$m.
We report what happens when every observable this programme can derive --- without passing through a
compactification that has never been constructed --- is written down, frozen by a git tag together
with its decision rule \emph{before} the verdict is computed, and then confronted with pinned data.
There are exactly four such observables. (P1) $s$ as the radius of one large extra dimension is
\textbf{excluded} by the E\"ot-Wash torsion-balance bound ($<30\,\mu$m) and fails the neutron-star bound
($<44\,\mu$m); the programme's own T-duality fixes the order-one factor that could have rescued it.
(C-A) Vacuum energy saturating the Cohen--Kaplan--Nelson bound with infrared length $c/H$ forces a
deceleration parameter $q=+1/2$ at every epoch, for every value of its only parameter, and is
\textbf{excluded} by the observed acceleration. (C-B) Pre-big-bang relic gravitons at the programme's
string scale end at $\approx6\times10^{-5}$~Hz with a peak $\Omega_{\rm GW}\approx10^{-65}$, more than fifty
orders of magnitude below LISA's design sensitivity: \textbf{not falsifiable in practice}. (C-C) Cosmic
F-strings have $G\mu\sim10^{-62}$: likewise out of reach. The premise shared by C-B and C-C --- a meV
string scale --- was already excluded by LHC resonance searches. The common cause of all four negative
verdicts is the identification $\alpha'=\ell_P c/H_0$ itself. Every derivation and every numerical
comparison is kernel-checked in Lean~4 (\texttt{DualScaleCosmology}, 50 theorems; repository 621 theorems,
0 failing, release \texttt{v3.16.1}); the physical identifications are stated as conjectures and the
experimental bounds as literature, each pinned to a source file and line range. We record the two freezes
made after the relevant bounds were known as retrodictions, and the one made before its comparison
document was pinned as a forward prediction. The mathematics of the programme (moonshine, dyons, lattices)
is unaffected and is reported separately; the lesson is methodological: pre-registration turned an
``order-one tension'' into a recorded exclusion and located the single assumption responsible.
\end{abstract}

\tableofcontents

%%%%%%%%%%%%%%%%%%%%%%%%%%%%%%%%%%%%%%%%%%%%%%%%%%%%%%%%%%%%%%%%%%%%%%%%%%%%%%%
\section{Introduction}
\label{sec:intro}
%%%%%%%%%%%%%%%%%%%%%%%%%%%%%%%%%%%%%%%%%%%%%%%%%%%%%%%%%%%%%%%%%%%%%%%%%%%%%%%

Papers~7--9 of this series~\cite{Callens2026Paper7,Callens2026Paper8,Callens2026Paper9} built a
Lean~4 companion to a ``dual-scale'' reading of string theory on $K3\times T^2$ and tested its literal
form against the Cohen--Kaplan--Nelson (CKN) bound~\cite{CKN1999}, which it fails by a factor of at least
$10^{30}$. Paper~9 left one reading open: the self-dual length $s=\sqrt{\ell_P\,c/H_0}\approx47\,\mu$m as
the radius of an extra dimension, ``disfavoured by an order-one factor, not excluded''. An order-one
tension is exactly where a theory can be quietly rescued by a convention chosen after the data are seen.
This paper closes that door. It asks two questions, in this order:
\begin{enumerate}
\item What in the programme yields a number \emph{with units}, by a derivation that does not pass through an
  unconstructed compactification (flux quanta, brane wrapping, moduli stabilisation)?
\item For each such number: what does a prediction frozen \emph{before} the verdict --- with its decision rule
  fixed in advance --- say, and what do pinned data say about it?
\end{enumerate}

\begin{headlinebox}
Four observables are derivable; all four were frozen and confronted. Two are \textbf{excluded} (P1, C-A);
two are \textbf{out of reach} of any foreseeable experiment (C-B, C-C), and their shared premise is
itself excluded. The common cause is the dual-scale identification $\alpha'=\ell_P c/H_0$. The testable
forms of the hypothesis do not survive existing data; the programme's Tier~A mathematics does.
\end{headlinebox}

\paragraph{Companion material.} Lean sources: \texttt{DualScaleCosmology/\allowbreak Stream6Verdict.lean},
\texttt{Stream7CA.lean}, \texttt{Stream7CB.lean}; they build on \texttt{SelfDualCutoff.lean},
\texttt{DarkEnergyScale.lean}, \texttt{CosmicString.lean}, \texttt{CKNInstance.lean}. Frozen predictions:
\texttt{docs/\allowbreak STREAM6\_PREDICTION\_P1.md}, \texttt{docs/\allowbreak STREAM7\_PREDICTION\_CA.md},
\texttt{docs/\allowbreak STREAM7\_PREDICTION\_CB.md}. Inventory: \texttt{docs/\allowbreak STREAM7\_HYPOTHESIS\_INVENTORY.md}.

%%%%%%%%%%%%%%%%%%%%%%%%%%%%%%%%%%%%%%%%%%%%%%%%%%%%%%%%%%%%%%%%%%%%%%%%%%%%%%%
\section{Epistemic tiers and the pre-registration protocol}
\label{sec:protocol}
%%%%%%%%%%%%%%%%%%%%%%%%%%%%%%%%%%%%%%%%%%%%%%%%%%%%%%%%%%%%%%%%%%%%%%%%%%%%%%%

Every claim carries one of three labels. \tierA: a proposition checked by the Lean~4 kernel (toolchain
\texttt{v4.33.1}, Mathlib), with no \texttt{sorry} and with axioms among \texttt{propext},
\texttt{Classical.choice}, \texttt{Quot.sound} (audited per theorem). \tierL: a statement taken from the
literature, pinned to a hash-registered file in \texttt{papers/foundations/} and a line range. \tierC: a
physical identification or hypothesis --- never a theorem. A Tier~A theorem about numbers with Tier~L
inputs certifies the arithmetic, not the inputs; a verdict inherits the weakest tier of its premises.

\begin{protocolbox}
\begin{enumerate}
\item \textbf{Derive before comparing.} The observable and its numerical value are derived from the
  programme's stated assumptions, each labelled measured (M), derived in Lean (D) or assumed (A). A non-empty
  (A) list makes the prediction Tier~C, stated as such.
\item \textbf{Freeze by tag.} The statement, the inputs and a \emph{decision rule} with outcome categories are
  committed and tagged (\texttt{stream6-p1-frozen}, \texttt{stream7-ca-frozen}, \texttt{stream7-cb-frozen})
  before the verdict is computed. Any later change is a new prediction, not an edit.
\item \textbf{Disclose prior knowledge.} If the relevant bound was already known to the authors, the freeze
  is recorded as a \emph{retrodiction}; only a freeze made before its comparison document was pinned is a
  \emph{forward} prediction.
\item \textbf{Compute the verdict mechanically} in Lean against the pinned bounds, and record it whatever it
  is --- including ``not falsifiable in practice'', which by this protocol is not a test.
\item \textbf{No post-hoc conventions.} An order-one factor may not be chosen after the bounds are read; it
  must be derived, or the prediction is taken with the convention stated in the freeze.
\end{enumerate}
\end{protocolbox}

Two of the four freezes (P1, C-A) are retrodictions: the bounds, and in C-A's case the outcome, were known
in advance, and the frozen documents say so. The C-B freeze preceded the pinning of the detector document it
was compared with.

%%%%%%%%%%%%%%%%%%%%%%%%%%%%%%%%%%%%%%%%%%%%%%%%%%%%%%%%%%%%%%%%%%%%%%%%%%%%%%%
\section{The programme's one physical identification}
\label{sec:identification}
%%%%%%%%%%%%%%%%%%%%%%%%%%%%%%%%%%%%%%%%%%%%%%%%%%%%%%%%%%%%%%%%%%%%%%%%%%%%%%%

The dual-scale programme pairs a micro length $\ell_{\rm micro}=\ell_P$ with a macro length
$\ell_{\rm macro}=c/H_0$ as a T-dual pair with $\alpha'=\ell_{\rm micro}\ell_{\rm macro}$ \tierC. Its
self-dual length is $s=\sqrt{\alpha'}$. With $\ell_P=1.616255\times10^{-35}$~m (CODATA~2018) and
$c/H_0=1.3672\times10^{26}$~m ($H_0=67.66$~km/s/Mpc, Planck~2018~VI~\cite{Planck2018VI}) \tierL,
$s^2\approx2.2097\times10^{-9}\,\mathrm{m}^2$, $s\approx47.0\,\mu$m, and $(40\,\mu\mathrm{m})^2\le s^2\le
(70\,\mu\mathrm{m})^2$ (\lean{SelfDualCutoff.selfDual\_length\_sq\_bracket}) \tierA.

Three facts from paper~9 frame what follows.
\begin{itemize}
\item \textbf{CKN.} With $L=\ell_{\rm macro}$, $M=1/\ell_{\rm micro}$ and UV length $\ell_{\rm UV}$, the corrected
  CKN bound holds iff $\ell_{\rm micro}\ell_{\rm macro}\le\ell_{\rm UV}^2$
  (\lean{SelfDualCutoff.ckn\_iff\_uvLength\_ge\_selfDual}) \tierA; the literal hypothesis (Planck-scale cutoff with
  Hubble IR) fails by at least $10^{30}$ (\lean{CKNInstance.planckEnergy\_gt\_ckn\_horizon\_cutoff}) \tierA.
\item \textbf{Regge reading excluded.} If $\alpha'=s^2$ is the Regge slope, string resonances sit near
  $\hbar c/s\approx4$~meV; CMS excludes them below 7.9~TeV in low-string-scale models~\cite{CMS2020} \tierL,
  and $s^2$ exceeds the CMS ceiling $(\hbar c/7.9\,\mathrm{TeV})^2$ by at least $10^{30}$
  (\lean{CosmicString.selfDual\_alphaPrime\_exceeds\_cms\_ceiling}) \tierA.
\item \textbf{Dark-energy identity.} With $\rho_\Lambda=\Omega_\Lambda\,3H_0^2/(8\pi G)$,
  $\rho_\Lambda^{-1}=(8\pi/3\Omega_\Lambda)(\ell_P L_H)^2$ (\lean{DarkEnergyScale.rhoLambda\_inv\_eq}) \tierA: the
  self-dual length \emph{is} the dark-energy length up to $(8\pi/3\Omega_\Lambda)^{1/4}$ --- an identity, not a
  corroboration.
\end{itemize}

%%%%%%%%%%%%%%%%%%%%%%%%%%%%%%%%%%%%%%%%%%%%%%%%%%%%%%%%%%%%%%%%%%%%%%%%%%%%%%%
\section{Inventory: what can be derived?}
\label{sec:inventory}
%%%%%%%%%%%%%%%%%%%%%%%%%%%%%%%%%%%%%%%%%%%%%%%%%%%%%%%%%%%%%%%%%%%%%%%%%%%%%%%

\Cref{tab:inventory} lists every output of the programme's five formal streams and asks whether it is a
number with units derivable without an unconstructed compactification.

\begin{table}[h]
\centering\small
\begin{tabular}{p{1.4cm}p{6.2cm}p{6.6cm}}
\toprule
Stream & Output & Observable? \\
\midrule
1--2 & lattices, $O(d,d;\mathbb Z)$, double field theory, the dual-scale bound, flux/tadpole arithmetic & no --- no number with units \\
3 & self-dual length $s$ & yes --- \textbf{P1} (\Cref{sec:p1}) \\
3 & CKN saturation with infrared length $c/H(t)$ & yes, given a Tier~C identification --- \textbf{C-A} (\Cref{sec:ca}) \\
3 & scale-factor duality $a\leftrightarrow1/a$ & only with dilaton-gravity dynamics --- \textbf{C-B} (\Cref{sec:cb}) \\
3 & cosmic F-string tension from $\alpha'=s^2$ & yes --- \textbf{C-C} (\Cref{sec:cc}) \\
4--5 & moonshine, umbral forms, dyon counts on $K3\times T^2$ & no --- dimensionless, and properties of $\mathcal N=4$ string theory \\
\bottomrule
\end{tabular}
\caption{Inventory of derivable observables (\texttt{docs/STREAM7\_HYPOTHESIS\_INVENTORY.md}).}
\label{tab:inventory}
\end{table}

Streams 4--5 deserve a word. Type~IIA string theory on $K3\times T^2$ gives an $\mathcal N=4$ theory in four
dimensions (Aspinwall, l.~2743~\cite{Aspinwall1996}) \tierL; four-dimensional $\mathcal N=4$ supersymmetry is
non-chiral (l.~2715) and admits no chiral matter (standard background, \tierL). The moonshine and dyon
results~\cite{Callens2026Paper10} are therefore mathematics of a supersymmetric string compactification, not
a model of our universe, and yield no observable here.

%%%%%%%%%%%%%%%%%%%%%%%%%%%%%%%%%%%%%%%%%%%%%%%%%%%%%%%%%%%%%%%%%%%%%%%%%%%%%%%
\section{P1: the self-dual length as an extra dimension}
\label{sec:p1}
%%%%%%%%%%%%%%%%%%%%%%%%%%%%%%%%%%%%%%%%%%%%%%%%%%%%%%%%%%%%%%%%%%%%%%%%%%%%%%%

\paragraph{Frozen statement} (tag \texttt{stream6-p1-frozen}, commit \texttt{53f15b7}; retrodiction,
disclosed). One large extra dimension of radius $R=s$, with the Kaluza--Klein convention $m\sim1/R$ used by
Montero--Vafa--Valenzuela (MVV, $m^{-1}\sim l$, l.~327~\cite{MVV2023}) \tierC.

\paragraph{Tests fixed in advance} \tierL:
\begin{itemize}
\item T1: E\"ot-Wash 2020, ``the largest extra dimension must have a toroidal radius less than
  $30\,\mu$m'' (l.~289~\cite{LeeAdelberger2020});
\item T2: any gravitational-strength Yukawa interaction has range $\lambda<38.6\,\mu$m (l.~285), read with
  $\lambda=R$ for the lowest KK mode (\tierC);
\item T3: heating of old neutron stars, one extra dimension, $l<44\,\mu$m (MVV l.~327).
\end{itemize}
Outcome rule: \emph{excluded} if T1 fails; \emph{in tension} if only T2 or T3 fails; \emph{consistent} otherwise.

\begin{theorem}[Verdict on P1] \tierA\ \label{thm:p1}
With $R^2=\ell_P\,c/H_0$ (\lean{Stream6Verdict.\allowbreak p1RadiusSq}):
$(30\,\mu\mathrm m)^2<R^2$, $(38.6\,\mu\mathrm m)^2<R^2$ and $(44\,\mu\mathrm m)^2<R^2$ (\lean{p1\_fails\_T1}, \lean{p1\_fails\_T2}, \lean{p1\_fails\_T3};
\lean{p1\_verdict\_excluded}). Margins: $R>1.5\times30\,\mu$m and $R>1.06\times44\,\mu$m (\lean{p1\_margin}).
Passing T1 would need $R=\kappa s$ with $\kappa<0.64$ (\lean{p1\_factor\_needed}: $\kappa=0.64$ fails,
$\kappa=0.63$ passes).
\end{theorem}
By the frozen rule, \textbf{P1 is excluded}. The factor $\kappa$ is recorded, not chosen: with $R=s/2\pi\approx
7.5\,\mu$m the number would fall inside MVV's ``dark dimension'' window $0.1$--$10\,\mu$m
(l.~427~\cite{MVV2023}), and choosing that convention after reading the bounds is exactly what the protocol forbids.

\paragraph{Can $\kappa$ be derived? Yes, and it is 1.} Within the programme's own T-duality $R\mapsto\alpha'/R$
with $\alpha'=s^2$:
\begin{theorem}[The self-dual radius is fixed] \tierA\ \label{thm:kappa}
For $R,s>0$: $R=s^2/R\iff R=s$ (\lean{p62\_selfdual\_fixed\_point}); at $R=s$ the Kaluza--Klein scale $1/R$ equals
the winding scale $R/\alpha'$ (\lean{p62\_towers\_coincide}).
\end{theorem}
With the standard normalisation $m_n=n/R$, the programme's radius is $\sqrt{\alpha'}=s$: $\kappa=1$. Taking both
circles of $T^2$ large makes things worse: MVV's two-dimension bound $l<1.6\times10^{-4}\,\mu$m (l.~329) is missed
by more than $10^5$ (\lean{p62\_two\_dims\_excluded}) \tierA.

%%%%%%%%%%%%%%%%%%%%%%%%%%%%%%%%%%%%%%%%%%%%%%%%%%%%%%%%%%%%%%%%%%%%%%%%%%%%%%%
\section{C-A: holographic dark energy at the Hubble scale}
\label{sec:ca}
%%%%%%%%%%%%%%%%%%%%%%%%%%%%%%%%%%%%%%%%%%%%%%%%%%%%%%%%%%%%%%%%%%%%%%%%%%%%%%%

\paragraph{Frozen statement} (tag \texttt{stream7-ca-frozen}, commit \texttt{90217a3}; retrodiction with a
foreseeable outcome, disclosed: Hubble-scale holographic dark energy is classically objected to because it does
not accelerate, and the acceleration is observed). Assumptions \tierC: (A1) $\rho_\Lambda(t)=\varepsilon\,
3H(t)^2/(8\pi G)$ with constant $0<\varepsilon<1$ (CKN saturation with infrared length $c/H$); (A2) flat FRW,
$H^2=(8\pi G/3)(\rho_m+\rho_\Lambda)$, with matter separately conserved, $\rho_m\propto a^{-3}$.

\begin{theorem}[Derivation] \tierA\ \label{thm:ca}
Write $k=8\pi G/3$. (i) (A1)--(A2) force $\rho_m=(1-\varepsilon)H^2/k$ (\lean{Stream7CA.ca\_matter\_fraction}).
(ii) Hence $H(a)^2=C/a^3$ with $C=k\rho_{m0}/(1-\varepsilon)$ (\lean{ca\_hubble\_scaling}). (iii) For $E(a)=C/a^3$,
$E'(a)=-3C/a^4$ (\lean{hasDerivAt\_E}) and the deceleration parameter $q=-1-aE'(a)/(2E(a))$ equals $1/2$ at every
$a>0$ and for every $C>0$ (\lean{ca\_deceleration}).
\end{theorem}
The prediction has no free parameter: whatever $\varepsilon$, the universe decelerates with $q=+1/2$.

\paragraph{Test and verdict.} Rule fixed in advance: C-A passes iff $q_0\ge0$. With Planck~2018's
$\Omega_\Lambda=0.68885$ in flat $\Lambda$CDM (\lean{DarkEnergyScale.omegaLambda}) \tierL, $q_0=(1-\Omega_\Lambda)/2-
\Omega_\Lambda\approx-0.53<0$, and $q_0\neq1/2$ (\lean{ca\_verdict\_excluded}) \tierA. \textbf{C-A is excluded.} The
input $q_0$ is model-dependent (a $\Lambda$CDM fit), as recorded in the freeze; the qualitative fact tested ---
present acceleration --- does not rest on $\Lambda$CDM alone, but no model-independent source is pinned here.

%%%%%%%%%%%%%%%%%%%%%%%%%%%%%%%%%%%%%%%%%%%%%%%%%%%%%%%%%%%%%%%%%%%%%%%%%%%%%%%
\section{C-B: pre-big-bang relic gravitons at the programme's string scale}
\label{sec:cb}
%%%%%%%%%%%%%%%%%%%%%%%%%%%%%%%%%%%%%%%%%%%%%%%%%%%%%%%%%%%%%%%%%%%%%%%%%%%%%%%

\paragraph{Frozen statement} (tag \texttt{stream7-cb-frozen}, commit \texttt{60320a4}; frozen \emph{before} any
detector document was pinned). Assumptions \tierC: (B1) the programme's scale-factor duality is realised
dynamically as a pre-big-bang phase; (B2) $\alpha'=s^2$ fixes the string scale, so the transition curvature
$g_1=H_1/M_P\simeq M_s/M_P=\ell_P/s$.

\paragraph{Literature input} \tierL\ (Gasperini--Veneziano~\cite{GasperiniVeneziano2002}). For $a=|\eta|^\alpha$ the
relic spectrum is $\Omega\sim\omega^{3-2\nu}$, $\nu=|\alpha-1/2|$ (Table~2, ll.~4467--4485); the dilaton-driven
phase has $\alpha=1/2$, hence $\Omega\sim\omega^3$ (l.~4505) in the low band (5.15), ll.~5040--5053; end points
$\omega_1\simeq g_1^{1/2}\,10^{11}$~Hz and $\Omega(\omega_1)\simeq10^{-4}g_1^2$, with $g_1=H_1/M_P\simeq M_s/M_P$
((5.16), ll.~5058--5061).

\begin{theorem}[Derived numbers] \tierA\ \label{thm:cb}
The slope $3-2|\alpha-1/2|$ is $3$ at $\alpha=1/2$ and $0$ (flat) at $\alpha=-1$ (\lean{Stream7CB.cb\_slope}).
With $g_1^2=\ell_P/(c/H_0)$ (\lean{g1Sq}): $1.1\times10^{-61}\le g_1^2\le1.3\times10^{-61}$ (\lean{cb\_g1\_sq\_bracket});
$\omega_1<10^{-4}$~Hz, stated as $\omega_1^4=g_1^2\cdot10^{44}<10^{-16}$ (\lean{cb\_peak\_below\_lisa\_band}); the
peak $\Omega(\omega_1)\simeq10^{-4}g_1^2$ (\lean{omegaPeak}) is below $10^{-64}$ (\lean{cb\_peak\_amplitude\_unreachable}).
\end{theorem}
So $\omega_1\approx6\times10^{-5}$~Hz and $\Omega_{\rm GW}(\omega_1)\approx1.2\times10^{-65}$; the only free element,
the break $\omega_s\le\omega_1$, can only lower the spectrum.

\paragraph{Test and verdict.} Rule fixed in advance: testable iff the prediction reaches the design sensitivity of
the most sensitive detector whose design document is pinned after the freeze. The LISA mission proposal requires
the ability to measure $\Omega=1.3\times10^{-11}(f/10^{-4}\,\mathrm{Hz})^{-1}$ for $0.1<f<2$~mHz (OR7.2,
ll.~692--699~\cite{LISA2017}) \tierL, i.e.\ at best $\approx6.5\times10^{-13}$ at 2~mHz (\lean{omegaLisaBest}). The
predicted spectrum ends below LISA's band and its peak, even multiplied by $10^{50}$, stays below LISA's best
(\lean{cb\_not\_testable}, \lean{cb\_peak\_amplitude\_unreachable}) \tierA. \textbf{C-B is not falsifiable in
practice} --- neither confirmed nor excluded by its rule.

\paragraph{Stronger.} Premise (B2) is the Regge reading of $\alpha'=s^2$, a meV string scale, which is already
excluded (\Cref{sec:identification}; \lean{CosmicString.selfDual\_alphaPrime\_exceeds\_cms\_ceiling}). C-B fails at the
level of its assumption, not only for lack of sensitivity.

%%%%%%%%%%%%%%%%%%%%%%%%%%%%%%%%%%%%%%%%%%%%%%%%%%%%%%%%%%%%%%%%%%%%%%%%%%%%%%%
\section{C-C: cosmic F-strings}
\label{sec:cc}
%%%%%%%%%%%%%%%%%%%%%%%%%%%%%%%%%%%%%%%%%%%%%%%%%%%%%%%%%%%%%%%%%%%%%%%%%%%%%%%

An unwarped cosmic F-string has $G\mu/c^2=\ell_P^2/(2\pi\alpha')$ (Copeland--Myers--Polchinski~\cite{CMP2004})
\tierL. For $\alpha'=\ell_P L$ this is $\ell_P/(2\pi L)$ (\lean{CosmicString.fStringGmu\_selfDual}) \tierA; with
$L=c/H_0$, $G\mu\approx1.9\times10^{-62}$, and in any case $G\mu\le10^{-61}$, more than fifty orders below the lower
end $10^{-11}$ of the brane-inflation window (\lean{CosmicString.selfDual\_Gmu\_below\_cmp\_window}) \tierA. A
prediction that no experiment can reach is not a test; C-C is recorded as out of reach (paper~9 deliberately does
not count the Planck $G\mu$ bound, passed by 55 orders, as evidence). It shares premise (B2) with C-B.

%%%%%%%%%%%%%%%%%%%%%%%%%%%%%%%%%%%%%%%%%%%%%%%%%%%%%%%%%%%%%%%%%%%%%%%%%%%%%%%
\section{The common cause}
\label{sec:cause}
%%%%%%%%%%%%%%%%%%%%%%%%%%%%%%%%%%%%%%%%%%%%%%%%%%%%%%%%%%%%%%%%%%%%%%%%%%%%%%%

\begin{table}[h]
\centering\small
\begin{tabular}{llp{4.2cm}p{4.6cm}}
\toprule
Observable & Freeze & Verdict & Where $\alpha'=\ell_Pc/H_0$ enters \\
\midrule
P1 & retrodiction & \textbf{excluded} (T1--T3) & $R=\sqrt{\alpha'}$, fixed by T-duality \\
C-A & retrodiction & \textbf{excluded} ($q=+1/2$ vs $q_0<0$) & IR length $c/H$ = $\ell_{\rm macro}$ \\
C-B & forward & not falsifiable in practice; premise excluded & $M_s\sim1/\sqrt{\alpha'}\sim4$~meV \\
C-C & (paper 9) & out of reach; premise excluded & $G\mu=\ell_P^2/(2\pi\alpha')$ \\
\bottomrule
\end{tabular}
\caption{The four verdicts and the role of the one identification.}
\label{tab:verdicts}
\end{table}

Each verdict traces back to the identification $\alpha'=\ell_{\rm micro}\ell_{\rm macro}=\ell_P c/H_0$
(\Cref{tab:verdicts}). Read as a length, it lands at tens of microns, where gravity has been tested and
it conflicts with data (P1). Read as an energy density at the Hubble scale, it fixes the dynamics
with no room for acceleration (C-A). Read as a string scale, it is a meV scale: excluded by colliders and, as a
by-product, hiding every stringy relic below any sensitivity (C-B, C-C). Together with the CKN result, the
dual-scale hypothesis in all of its testable forms does not survive the existing data.

%%%%%%%%%%%%%%%%%%%%%%%%%%%%%%%%%%%%%%%%%%%%%%%%%%%%%%%%%%%%%%%%%%%%%%%%%%%%%%%
\section{What survives}
\label{sec:survives}
%%%%%%%%%%%%%%%%%%%%%%%%%%%%%%%%%%%%%%%%%%%%%%%%%%%%%%%%%%%%%%%%%%%%%%%%%%%%%%%

The verdicts concern physical identifications (Tier~C). The Tier~A mathematics of the programme is untouched:
lattices, $O(d,d;\mathbb Z)$ and the dual-scale bound $\operatorname{tr}G+\operatorname{tr}G^{-1}\ge2d$ with its unique
minimiser (\lean{dualScale\_eq\_iff}, paper~8~\cite{Callens2026Paper8}); Mathieu moonshine computed at all 26 classes
and levels 0--9 (\lean{moonshine\_modules\_decompose}), the shadow's $24$ (\lean{shadow\_coeff\_eq\_perm\_trace}), and the
single-centred dyon counts on $K3\times T^2$ given by Hurwitz class numbers (\lean{immortal\_m1}), all in
paper~10~\cite{Callens2026Paper10}. The same twining (``forger's'') test that the number $24=\chi(K3)$ passes is
failed by paper~7's ``27720 lock'' at every non-identity class (\lean{ratio\_fails\_at\_every\_class}); paper~7's
Revision~5 records this.

\begin{scopebox}
That any of the identifications (P1, A1--A2, B1--B2) is part of string theory (Tier~C); that the Tier~L inputs are
correct beyond their sources; a model-independent measurement of $q_0$ (not pinned); any realistic compactification.
Two of four freezes are retrodictions, as disclosed. ``Not falsifiable in practice'' is not a confirmation.
\end{scopebox}

%%%%%%%%%%%%%%%%%%%%%%%%%%%%%%%%%%%%%%%%%%%%%%%%%%%%%%%%%%%%%%%%%%%%%%%%%%%%%%%
\section{Lessons for pre-registered theory testing}
\label{sec:lessons}
%%%%%%%%%%%%%%%%%%%%%%%%%%%%%%%%%%%%%%%%%%%%%%%%%%%%%%%%%%%%%%%%%%%%%%%%%%%%%%%

\begin{enumerate}
\item \textbf{Inventory before hypothesis.} Asking first which observables are derivable at all prevented a search
  for a convenient one; the answer (four) bounded the whole exercise.
\item \textbf{Freeze the order-one factors.} Paper~9's ``disfavoured by $O(1)$'' became an exclusion once the factor was
  either fixed in the freeze or derived (\Cref{thm:kappa}). An unfixed factor is where post-hoc rescue lives.
\item \textbf{Retrodiction is still worth recording.} Even with a known outcome, a frozen, kernel-checked derivation
  (C-A) turns an informal objection into a checkable one and removes it from future debate.
\item \textbf{Separate ``not excluded'' from ``not testable''.} C-B and C-C are neither confirmed nor excluded by
  their rules; the protocol says so explicitly rather than counting consistency with an unreachable bound as
  support.
\item \textbf{Name quantities in verdict theorems.} A review of release \texttt{v3.16.0} found two C-B statements whose
  second conjunct was a fact about powers of ten; they were restated with named quantities (\lean{omegaPeak},
  \lean{omegaLisaBest}), with the statement lock recording exactly those two changes.
\end{enumerate}

%%%%%%%%%%%%%%%%%%%%%%%%%%%%%%%%%%%%%%%%%%%%%%%%%%%%%%%%%%%%%%%%%%%%%%%%%%%%%%%
\section{Open directions}
\label{sec:open}
%%%%%%%%%%%%%%%%%%%%%%%%%%%%%%%%%%%%%%%%%%%%%%%%%%%%%%%%%%%%%%%%%%%%%%%%%%%%%%%

\begin{enumerate}
\item \textbf{Drop or replace $\alpha'=\ell_Pc/H_0$.} Any new hypothesis worth testing must do without this
  identification; its observables must again be derived and frozen before comparison.
\item \textbf{A chiral geometry.} $K3\times T^2$ gives $\mathcal N=4$ in four dimensions; a realistic model needs a
  geometry with a chiral spectrum --- for instance a $K3$-fibred Calabi--Yau threefold --- together with an explicit
  construction (fluxes, moduli stabilisation). This is new research, not a patch.
\item \textbf{Model-independent inputs.} Pinning a model-independent determination of the present acceleration
  would lift C-A's verdict from a $\Lambda$CDM-conditional one.
\item \textbf{Reusable protocol.} The freeze-by-tag protocol and the Lean verdict files are generic and could be
  reused by other programmes that make order-one-sensitive claims.
\end{enumerate}

%%%%%%%%%%%%%%%%%%%%%%%%%%%%%%%%%%%%%%%%%%%%%%%%%%%%%%%%%%%%%%%%%%%%%%%%%%%%%%%
\section{Conclusion}
\label{sec:conclusion}
%%%%%%%%%%%%%%%%%%%%%%%%%%%%%%%%%%%%%%%%%%%%%%%%%%%%%%%%%%%%%%%%%%%%%%%%%%%%%%%

A hypothesis that links the Planck length and the Hubble radius through string T-duality yields exactly four
observables without an unconstructed compactification. Frozen before their verdicts and checked in Lean against
pinned data, two are excluded and two are out of reach with an excluded premise. The single cause is the
identification $\alpha'=\ell_Pc/H_0$. We regard this as a clean negative result: the programme's testable physical
content does not survive, its mathematics does, and the protocol that established this is reusable.

%%%%%%%%%%%%%%%%%%%%%%%%%%%%%%%%%%%%%%%%%%%%%%%%%%%%%%%%%%%%%%%%%%%%%%%%%%%%%%%
\appendix
\section{Reproducibility}
\label{app:repro}
%%%%%%%%%%%%%%%%%%%%%%%%%%%%%%%%%%%%%%%%%%%%%%%%%%%%%%%%%%%%%%%%%%%%%%%%%%%%%%%

\paragraph{Build and audit.}
{\footnotesize\ttfamily lake build DualScaleCosmology}\\
{\footnotesize\ttfamily python3 tools/axiom\_audit.py DualScaleCosmology}\\
reports 50 theorems, 0 failing; across the ten libraries of the repository, 621 theorems, 0 failing (release
\texttt{v3.16.1}). Axioms used: at most \texttt{propext}, \texttt{Classical.choice}, \texttt{Quot.sound}.

\paragraph{Statement lock.}
{\footnotesize\ttfamily python3 tools/statement\_lock.py --check \$(find DualScaleCosmology -name '*.lean')}\\
locks 68 declarations in 10 files, including \texttt{Stream6Verdict.lean}, \texttt{Stream7CA.lean} and
\texttt{Stream7CB.lean}.

\paragraph{Freezes.} \texttt{git show stream6-p1-frozen}, \texttt{stream7-ca-frozen}, \texttt{stream7-cb-frozen}
(commits \texttt{53f15b7}, \texttt{90217a3}, \texttt{60320a4}) display the frozen documents; the verdict files were
committed afterwards.

\paragraph{Sources pinned} (hash-registered in \texttt{papers/foundations/MANIFEST.md}): \texttt{2002.11761},
\texttt{2205.12293}, \texttt{hep-th/0207130}, \texttt{1702.00786}, \texttt{hep-th/0312067}, \texttt{1911.03947},
\texttt{hep-th/9803132}, \texttt{hep-th/9611137}; CODATA~2018 and Planck~2018~VI constants via \texttt{astropy}.

\begin{thebibliography}{99}
\bibitem{Callens2026Paper7} X.~Callens and the SocrateAI Agora Collaboration, ``The Dual-Scale String Theory: Mechanized Foundations, Singularity Resolution, Mathieu Moonshine, and a Zero-Free-Parameter Cosmological Conjecture on $K3\times T^2$,'' Revision~5 (2026).
\bibitem{Callens2026Paper8} X.~Callens and the SocrateAI Agora Collaboration, ``Lattices, T-Duality, and Double Field Theory on $K3\times T^2$: A Lean 4 Companion Formalization'' (2026).
\bibitem{Callens2026Paper9} X.~Callens and the SocrateAI Agora Collaboration, ``Testing a Micro/Macro Dual-Scale Hypothesis on $K3\times T^2$: Scale-Factor Duality, the Cohen--Kaplan--Nelson Bound, and the Dark-Energy Length in Lean~4'' (2026).
\bibitem{Callens2026Paper10} X.~Callens and the SocrateAI Agora Collaboration, ``Mathieu Moonshine Computed, the Double-Scaled Little String Bridge, and Dyons on $K3\times T^2$ in Lean~4'' (2026).
\bibitem{CKN1999} A.~G.~Cohen, D.~B.~Kaplan, and A.~E.~Nelson, ``Effective Field Theory, Black Holes, and the Cosmological Constant,'' \textit{Phys. Rev. Lett.} 82 (1999) 4971, \textit{arXiv:hep-th/9803132}.
\bibitem{Planck2018VI} Planck Collaboration, ``Planck 2018 results. VI. Cosmological parameters,'' \textit{Astron. Astrophys.} 641 (2020) A6, \textit{arXiv:1807.06209}.
\bibitem{CMS2020} CMS Collaboration, ``Search for high mass dijet resonances with a new background prediction method in proton-proton collisions at $\sqrt{s}=13$~TeV,'' \textit{JHEP} 05 (2020) 033, \textit{arXiv:1911.03947}.
\bibitem{LeeAdelberger2020} J.~G.~Lee, E.~G.~Adelberger, T.~S.~Cook, S.~M.~Fleischer, and B.~R.~Heckel, ``New Test of the Gravitational $1/r^2$ Law at Separations down to $52\,\mu$m,'' \textit{Phys. Rev. Lett.} 124 (2020) 101101, \textit{arXiv:2002.11761}.
\bibitem{MVV2023} M.~Montero, C.~Vafa, and I.~Valenzuela, ``The Dark Dimension and the Swampland,'' \textit{JHEP} 02 (2023) 022, \textit{arXiv:2205.12293}.
\bibitem{GasperiniVeneziano2002} M.~Gasperini and G.~Veneziano, ``The Pre-Big Bang Scenario in String Cosmology,'' \textit{Phys. Rept.} 373 (2003) 1, \textit{arXiv:hep-th/0207130}.
\bibitem{LISA2017} P.~Amaro-Seoane et al., ``Laser Interferometer Space Antenna,'' \textit{arXiv:1702.00786} (2017).
\bibitem{CMP2004} E.~J.~Copeland, R.~C.~Myers, and J.~Polchinski, ``Cosmic F- and D-strings,'' \textit{JHEP} 06 (2004) 013, \textit{arXiv:hep-th/0312067}.
\bibitem{Aspinwall1996} P.~S.~Aspinwall, ``K3 Surfaces and String Duality,'' \textit{arXiv:hep-th/9611137} (1996).
\end{thebibliography}

\end{document}
